%% Section 2 -- Temporal Communication Fragmentation
%% Ported from docs/DRAFT_PAPER.md Section 2. British spelling is deliberate.

\section{Temporal Communication Fragmentation}
\label{sec:tcf}

\subsection{Definition}
\label{sec:tcf:definition}

\textbf{Temporal Communication Fragmentation} is a state in which a user's
communication history for a single logical account becomes fragmented across
two or more device-bound application instances that persist independently, and
for which the platform provides no operation yielding their union. It arises
from primary-device unavailability, continued account use on a secondary
device during that window, and primary-device recovery.

The definition requires the absence of a union-yielding operation and nothing
further. How a given platform declines to provide one is a separate, empirical
matter. On the platform examined here, every path that brings the archives back
together, restore above all, substitutes one state for another rather than
merging them, with no user-facing warning that a competing archive exists and
no attempt at a diff or a reconciliation. That is an observation about one
documented restore path, not a property TCF requires.

We assert the observable restore behaviour and nothing about the backend that
produces it. A single-slot, last-write-wins cloud backup destination would
explain what users see, but we have found no vendor documentation, citation, or
controlled experiment establishing that architecture, and no argument in this
paper depends on it.

%% Figure 1 -- TCF timeline. Drawn from Section 2.2 only: the three necessary
%% conditions, the t_f / t_r landmarks, and the F_1 / F_2 extents with delta
%% and epsilon. Nothing here is asserted that Section 2.2 does not already
%% state. delta and epsilon are drawn strictly positive because that is the
%% case the prose singles out as the reason deduplication is needed at all
%% (Sections 2.2 and 3.4); both are defined as merely non-negative, and the
%% caption says so rather than letting the drawing imply otherwise.
%%
%% NB (2026-08-26, ground-problem alignment): C1 and C2 are events and sit on
%% the timeline; C3 is a standing property of the platform and must NOT be
%% drawn at a point. It spans from t_f onwards. The restore is drawn dashed
%% and labelled as the manifest case only -- conditions 1-3 are already
%% satisfied without it. Reverting either of those would reintroduce the
%% overfitting this pass removed.
\begin{figure}[htbp]
  \centering
  \begin{tikzpicture}[
      x=1mm, y=1mm,
      >={Stealth[length=2mm]},
      frag/.style={rounded corners=1pt, draw, thick},
      lbl/.style={font=\small},
      tick/.style={draw, thick},
      cond/.style={font=\scriptsize\itshape, align=center},
    ]

    %% landmark positions
    \def\tf{34}   % device failure / bifurcation
    \def\tr{86}   % primary device recovery
    \def\trest{110} % restore, if it is ever performed
    \def\dlt{10}  % delta : pre-existing history on the secondary
    \def\eps{8}   % epsilon: lag before the account returns to the primary

    %% ---- C3 : a standing property, not an event -------------------------
    \draw[decorate, decoration={brace, amplitude=4pt}] (\tf,41) -- (126,41);
    \node[cond, above=1.5mm] at ({(\tf+126)/2},41)
      {\textbf{C3} Absent reconciliation: no operation yields
       $F_1 \cup F_2$\\(holds from $t_f$ onwards, whether or not
       a restore is ever performed)};

    %% ---- time axis -------------------------------------------------------
    \draw[->, thick] (0,0) -- (128,0) node[right, lbl] {time};

    \foreach \x in {\tf, \tr} {
      \draw[tick] (\x,-2) -- (\x,2);
    }
    \node[lbl, below=1mm] at (\tf,-2) {$t_f$};
    \node[lbl, below=1mm] at (\tr,-2) {$t_r$};

    %% ---- F_1 : primary -- paused across the window ------------------------
    \draw[frag, fill=black!8] (2,10) rectangle (\tf,18);
    \draw[frag, fill=black!8] (\tr,10) rectangle (126,18);
    \draw[dashed, gray] (\tf,14) -- (\tr,14);
    \node[lbl] at (18,14) {$F_1$};
    \node[lbl] at (118,14) {$F_1$};
    \node[cond, gray] at ({(\tf+\tr)/2},14) {primary paused\\(no new history)};

    %% ---- F_2 : secondary -- active across the window ----------------------
    \draw[frag, fill=black!20] ({\tf-\dlt},22) rectangle ({\tr+\eps},30);
    \node[lbl] at ({(\tf+\tr)/2},26) {$F_2$};

    %% ---- the two overlap regions -----------------------------------------
    \begin{scope}[on background layer]
      \fill[pattern=north east lines, pattern color=black!45]
        ({\tf-\dlt},10) rectangle (\tf,30);
      \fill[pattern=north east lines, pattern color=black!45]
        (\tr,10) rectangle ({\tr+\eps},30);
    \end{scope}
    \draw[<->] ({\tf-\dlt},34) -- (\tf,34)
      node[midway, above, font=\scriptsize] {$\delta$};
    \draw[<->] (\tr,34) -- ({\tr+\eps},34)
      node[midway, above, font=\scriptsize] {$\epsilon$};

    %% ---- latent TCF : C1-C3 hold, both archives still readable -----------
    %% Single line, anchored inside the 0--10 gap. Do not re-anchor this above
    %% its arrow: it then rides up into the right-hand F_1 rectangle.
    \draw[<->, gray] (\tr,7.5) -- (\trest,7.5);
    \node[cond, gray] at ({(\tr+\trest)/2},4)
      {latent: both readable};

    %% ---- the restore, if performed : latent becomes manifest -------------
    \draw[very thick, dashed] (\trest,8) -- (\trest,32);

    %% ---- below-axis annotations, in three non-overlapping bands ----------
    %% band 1 (y -8 to -26): C1 at t_f and the manifest note at the restore;
    %% they do not overlap horizontally. Arrows start at -8, below the tick
    %% labels, so they do not strike through $t_f$ / $t_r$.
    \draw[->, thick] (\tf,-15) -- (\tf,-8);
    \node[cond, below=0.5mm] at (\tf,-15)
      {\textbf{C1} Bifurcation\\primary pauses,\\secondary active};

    \draw[->, thick] (\trest,-15) -- (\trest,-8);
    \node[cond, below=0.5mm] at (\trest,-15)
      {\emph{manifest} case only\\restore replaces:\\$F_2$ no longer surfaced};

    %% band 2 (y -28 down): C2, spanning t_f to t_r, clear of band 1.
    \draw[decorate, decoration={brace, amplitude=4pt, mirror}]
      (\tf,-28) -- (\tr,-28);
    \node[cond, below=2mm] at ({(\tf+\tr)/2},-30)
      {\textbf{C2} Independent persistence\\both branches accumulate,\\zero synchronisation};

  \end{tikzpicture}
  %% Caption describes what is drawn. The argument behind the drawing --- why
  %% C3 is a span rather than a point, why the restore is optional, why the
  %% overlap needs deduplication --- is in Section 2.2 and must not migrate
  %% back up here (peer review, Minor 13: a caption doing expository work).
  \caption{Temporal Communication Fragmentation on a timeline, with the three
  necessary conditions of Section~\ref{sec:tcf:conditions} marked. $F_1$ holds
  messages with $t(m) \le t_f$ or $t(m) \ge t_r$; $F_2$ holds messages in
  $[t_f - \delta,\, t_r + \epsilon]$. Hatching marks the two overlap regions.
  $\delta$ and $\epsilon$ are non-negative in the model and are drawn here as
  strictly positive so that the overlap is visible. C1 and C2 are drawn as
  points and C3 as a span; the restore is dashed because it is optional.}
  \label{fig:tcf_timeline}
\end{figure}

\subsection{Necessary conditions}
\label{sec:tcf:conditions}

TCF requires all three of the following to hold at once
(Figure~\ref{fig:tcf_timeline}).

\paragraph{Condition 1, Bifurcation.}
At failure time $t_f$, one authoritative timeline splits into two independently
updating branches. The primary is paused; the secondary is active.

\paragraph{Condition 2, Independent persistence.}
Both branches survive and accumulate new history for some duration with zero
synchronisation between them. This condition does most of the work in separating
TCF from the replication literature, and Section~\ref{sec:related:replication}
returns to it.

\paragraph{Condition 3, Absent reconciliation.} No platform-native
operation yields the union of the two archives. This condition is a property of
the platform's reunification semantics rather than an event that must occur. It
holds from the moment the branches exist, whether or not the user ever invokes
one of those paths, and it is what separates TCF from the
optimistic-replication systems of Section~\ref{sec:related:replication}, which
attempt reconciliation imperfectly rather than declining to attempt it.

Whether a platform's substitutive paths also overwrite silently, warn, or offer
a comparison bears on how much a user can recover unaided, not on whether the
condition holds. In the restore path examined here none of those mitigations is
present: restoration overwrites current chat state instead of merging into it,
surfaces no warning, and gives no indication that a second archive existed.
That is the observed behaviour of one platform, recorded because it is the
motivating case, and Condition 3 does not depend on it.

Stated most carefully, Condition 3 is a property of the reconstruction problem's
input rather than a claim about any company's feature list: the fragments carry
no shared reconciliation state, and no union-yielding operation is available to
the user holding them. This matters for what a changelog can do to this paper. If
a platform ships a merge affordance tomorrow, instances of TCF on that platform
cease to arise, in the same way that a fix removes instances of a defect class
without refuting its description. What the definition would lose is reach, not
correctness --- and the reconstruction problem for archives already fragmented
before such a fix would be untouched, because the fragments in hand would still
carry no shared reconciliation state. A definition whose truth conditions sat
inside a vendor's roadmap would be a poor research object; this one names a
configuration that either obtains for a given pair of archives or does not.

\paragraph{Latent and manifest fragmentation.} Conditions 1 to 3 are satisfied
while both archives remain individually readable on their own devices; what is
unavailable is any view spanning them. We call this the \emph{latent} case,
and it is the case the motivating incident of Section~\ref{sec:introduction}
exhibits. A restore makes the fragmentation \emph{manifest}: the replaced
archive thereafter survives only in whatever unsynced local form remains on
the device holding it, and the platform ceases to surface it at all. Manifest
fragmentation is the severe case of TCF, not its definition. The distinction
matters for how urgent reconstruction is and for what a user can still recover
unaided, but not for the reconstruction problem itself, because both cases
present the same input: independently preserved fragments carrying no shared
reconciliation state.

Formally (Section~\ref{sec:carta} develops this), device failure at $t_f$ and
recovery at $t_r$ split a ground-truth history $\mathcal{T}$ into $F_1$,
holding messages with $t(m) \le t_f$ or $t(m) \ge t_r$, and $F_2$, holding
messages in $[t_f - \delta,\, t_r + \epsilon]$. The parameters $\delta$ and
$\epsilon$ are non-negative, representing pre-existing history on the
secondary device and the lag before the account returns to the primary. Both
may be strictly positive, so $F_1$ and $F_2$ can overlap. That overlap is why
reconstruction needs deduplication rather than concatenation.

%% Added 2026-08-26 after review: the motivating incident is the delta = 0
%% boundary, and the paper previously asserted overlap as though it were
%% general without saying where the general case comes from. Do not delete
%% without also softening Section 1's "overlap in part" and Section 3.4's
%% "deduplication exists because of this overlap".
The boundary case $\delta = \epsilon = 0$ deserves explicit mention, because it
is the one the motivating incident of Section~\ref{sec:introduction} exhibits:
the secondary device was set up without restoring any prior history, so $F_2$
begins empty at $t_f$ and $F_1 \cap F_2 = \emptyset$. Reconstruction there
reduces to ordered concatenation, and none of the deduplication machinery of
Section~\ref{sec:algorithms} is exercised. We do not treat that case as the
general one. A secondary device is more often brought into service the way any
replacement device is, by restoring whatever backup exists, and a backup made
before $t_f$ carries pre-failure history that the primary also still holds ---
which is exactly $\delta > 0$. Symmetrically, an account does not always return
to the primary the moment it is repaired, and any use of the secondary after
$t_r$ produces $\epsilon > 0$. The overlapping configuration is therefore the
one a specification has to be written for, and the disjoint one is the easy
special case that falls out of it. Assuming disjointness and discovering
overlap later would mean a reconstruction that silently reported the overlap
region twice.

Where this paper calls a fragment \emph{internally complete}, it means
well-formed and gap-free with respect to its own accumulation window, which is
a property of each fragment considered separately. It is not a claim that
either fragment covers $\mathcal{T}$; by construction neither does, and that
is the whole difficulty. Internal completeness is an idealisation of the TCF
model rather than a property established of any real archive, and
Section~\ref{sec:discussion:incomplete} treats the case where it fails.

\subsection{Distinction from backup staleness}
\label{sec:tcf:staleness}

Backup staleness is the documented problem of infrequent snapshots losing data
when a device fails. One timeline exists, and the data from the unprotected
window was never captured. TCF has two timelines, each holding messages the other
lacks, both physically present. The reconstruction problem is one of retrieval
and reconciliation rather than of recovery from destruction.

\subsection{Position against neighbouring phenomena}
\label{sec:tcf:neighbours}

%% Table 1 -- Taxonomy of neighbouring data-loss and consistency phenomena.
%% Reproduced unchanged from docs/DRAFT_PAPER.md Section 2.4, which in turn
%% reproduces docs/PAPER_BLUEPRINT.md Section 4. \input from 02_tcf.tex.
%%
%% Note: the column headings and the TCF column use the source documents'
%% spelling ("Platform Behavior") rather than being Anglicised, because the
%% table is reproduced verbatim. See docs/DRAFT_PAPER.md drafting note 5.

\begin{table}[htbp]
  \centering
  \caption{Taxonomy of neighbouring data-loss and consistency phenomena.}
  \label{tab:taxonomy}
  \small
  \setlength{\tabcolsep}{4pt}
  \renewcommand{\arraystretch}{1.25}
  \begin{tabularx}{\textwidth}{@{}l >{\raggedright\arraybackslash}X >{\raggedright\arraybackslash}X >{\raggedright\arraybackslash}X >{\raggedright\arraybackslash}X@{}}
    \toprule
    \textbf{Dimension} &
      \textbf{Backup Staleness} &
      \textbf{Sync Conflict} &
      \textbf{Git 3-Way Merge} &
      \textbf{TCF (This Paper)} \\
    \midrule
    \textbf{Underlying Cause} &
      Infrequent snapshotting &
      Concurrent multi-client writes &
      Divergent commit histories &
      Device switch + independent persistence; no merge operation \\

    \textbf{Number of Timelines} &
      Exactly 1 &
      2+ concurrent &
      2 branches + Common Ancestor &
      2 independently persistent, potentially overlapping \\

    \textbf{Platform Behavior} &
      Restores stale snapshot &
      Rejection or branch fork &
      Automatic merge or conflict markers &
      \textbf{No operation yields the union; every reunification path replaces} \\

    \textbf{Physical Data State} &
      Missing (never recorded) &
      Present on server &
      Merged in working tree &
      \textbf{Intact on both devices; no combined view} \\

    \textbf{Resolution Target} &
      Data unrecoverable from the snapshot &
      Operational state restored &
      Codebase reconciliation &
      \textbf{Forensic reconstruction ($\mathcal{U}, \Omega$)} \\
    \bottomrule
  \end{tabularx}
\end{table}


Two further neighbours deserve separate treatment because they are close enough
to invite misreading, and both are already-named phenomena rather than gaps.

\paragraph{Sync conflict.}
Systems built for optimistic replication also allow replicas to diverge, and they
define conflict as a first-class condition. \citet{saito2005optimistic} survey
the area and describe divergence and conflict as the characteristic problems such
systems face. The difference is what happens on reunification: those platforms
attempt reconciliation, imperfectly and with human fallback, whereas the platform
behaviour in Condition 3 makes no attempt at all.

\paragraph{Device migration.}
Moving an account between handsets is routine, and the HCI literature on
multi-device use describes the resulting failure as information stranded on the
device the user is not currently holding~\citep{dearman2008othercomputer}.
Migration moves one continuous timeline. TCF requires two timelines that
accumulated independently and concurrently.

\subsection{Terminology: what TCF is not called}
\label{sec:tcf:terminology}

Three established terms sit close enough to ``temporal communication
fragmentation'' that a reader may reasonably assume the phenomenon is already
covered. Each denotes something else, and we set the boundaries explicitly rather
than rely on context.

\paragraph{Information fragmentation (personal information management).}
\citet{karger2006unification} define this as personal information split by the
very tools meant to manage it, with applications storing data ``in their own
particular locations and representations, inaccessible to other applications'';
the project-level variant is treated by Bergman and
colleagues~\citep{bergman2006project}. Later usage in the field routinely extends
the term across devices as well as applications, so the surface overlap with TCF
is real. Two things separate them. The unit differs: information fragmentation
concerns \emph{different items} bearing on one task or project, scattered across
stores, whereas TCF concerns \emph{one item stream}, a single conversation, split
into disjoint temporal segments. And the access status differs: fragmented
personal information is scattered but reachable, and the research goal is
unification across silos, whereas in TCF one intact archive is unreachable
through the platform that holds it. Karger and Jones do treat duplicate copies
drifting out of step, which is content divergence between copies a user made;
TCF's divergence is temporal, between copies the platform made and then rendered
mutually exclusive.

\paragraph{Fragmented communication (CSCW and workplace studies).}
In that literature the term describes one conversation spread across many
channels and tools. TCF is the opposite arrangement: one channel, one account,
one conversation, split by time rather than by venue.

\paragraph{Temporal fragmentation (time geography and ICT studies).}
\citet{hubers2008temporal} use this for the splitting of human \emph{activities}
into pieces performed at different times and places. The referent is an activity
pattern, not a data archive. We keep the name Temporal Communication
Fragmentation because the qualifier ``communication'' and the archival framing
distinguish the two, and because no prior work applies any of these three terms
to the configuration defined in Section~\ref{sec:tcf:conditions}. The collision
is nonetheless a live source of confusion, and it is the reason this subsection
exists.
